% Performance benchmark: full-feature document (realistic engineering report).
% Combines code blocks, math, TikZ, tables, figures, and bibliography.
% Expected compile time: <= 90 s (three passes with biber).
% !TEX program = xelatex
\documentclass[titlepage=formal,code=listings]{SUSTechHomework}

\title{Perf Benchmark: Full Features}
\coursecode{PERF}
\coursename{Full Feature Benchmark}
\semester{Spring 2026}
\instructor{Prof. Benchmark}
\reporttype{Performance Test}
\date{2026-04-01}
\addbibresource{example.bib}

\begin{document}

\clearauthors
\addauthor{Alice}{11000001}{alice@perf.test}{Department of Computer Science}{Shuren College}{Lead}
\addauthor{Bob}{11000002}{bob@perf.test}{Department of EEE}{Zhixin College}{Experiment}

\maketitle
\tableofcontents

% ------------------------------------------------------------
\section{Introduction}

This benchmark exercises every major subsystem of the template in a single
document to measure realistic compile time for a full engineering report.

\begin{sustechnote}[title=Benchmark Note]
This document is a performance regression fixture. It should compile within
the specified time limit on the reference machine.
\end{sustechnote}

% ------------------------------------------------------------
\section{Mathematics}

The response energy integral:
\[
  I = \int_{0}^{T} \e^{-t} \sin(\omega t)\,\diff t
  = \frac{\omega}{\omega^{2}+1}\bigl(1 - \e^{-T}\cos(\omega T)\bigr).
\]

Vector and matrix notation:
\begin{align}
  \mat{K}\vect{x} &= \vect{b}, \quad \vect{x} \in \mathbb{R}^{n}, \\
  \norm*{\vect{x}}_2 &\le \abs*{\alpha} + \qty{2.5}{\newton}, \\
  \vect{v}(t) &= \begin{bmatrix} x(t) \\ y(t) \\ z(t) \end{bmatrix}.
\end{align}

% ------------------------------------------------------------
\section{Code}

\begin{sustechcode}{Python}
import numpy as np

def simulate(A, B, x0, u, dt=0.01):
    """Discrete-time linear system simulation."""
    x = x0.copy()
    states = [x.copy()]
    for uk in u:
        x = A @ x + B * uk
        states.append(x.copy())
    return np.array(states)
\end{sustechcode}

\begin{sustechcode}{moderncpp}
#include <Eigen/Dense>
#include <vector>

Eigen::VectorXd simulate(
    Eigen::MatrixXd const& A,
    Eigen::VectorXd const& B,
    Eigen::VectorXd       x0,
    std::vector<double> const& u,
    double dt = 0.01)
{
    for (auto uk : u) {
        x0 = A * x0 + B * uk;
    }
    return x0;
}
\end{sustechcode}

\begin{sustechcode}{matlab}
function x = simulate(A, B, x0, u, dt)
    if nargin < 5, dt = 0.01; end
    x = x0;
    for k = 1:length(u)
        x = A * x + B * u(k);
    end
end
\end{sustechcode}

\begin{sustechdiff}
--- a/src/sim.py
+++ b/src/sim.py
@@ -3,6 +3,7 @@
 def simulate(A, B, x0, u):
-    states = []
+    states = [x0.copy()]
     x = x0.copy()
     for uk in u:
         x = A @ x + B * uk
+        states.append(x.copy())
     return states
\end{sustechdiff}

% ------------------------------------------------------------
\section{Tables}

\begin{table}[htbp]
  \centering
  \begin{tabular}{llll}
    \sustechheaderrow
    \sustechtableheader{Method} & \sustechtableheader{RMSE} &
    \sustechtableheader{MAE}  & \sustechtableheader{Time (s)} \\
    Baseline    & 0.312 & 0.245 & 1.2 \\
    Proposed    & 0.178 & 0.143 & 1.8 \\
    Ablation A  & 0.221 & 0.189 & 1.5 \\
    Ablation B  & 0.264 & 0.210 & 1.3 \\
  \end{tabular}
  \caption{Performance comparison across methods}
  \label{tab:perf-methods}
\end{table}

% ------------------------------------------------------------
\section{Figures}

\begin{sustechfiguregroup}
  \begin{sustechsubfigure}[0.48\linewidth][Baseline]
    \sustechplaceholdergraphic[0.95\linewidth][3.5cm]
  \end{sustechsubfigure}
  \sustechfloatsep
  \begin{sustechsubfigure}[0.48\linewidth][Proposed]
    \sustechplaceholdergraphic[0.95\linewidth][3.5cm]
  \end{sustechsubfigure}
  \caption{Visual comparison of baseline versus proposed method}
\end{sustechfiguregroup}

% ------------------------------------------------------------
\section{Diagram}

\begin{figure}[htbp]
  \centering
  \begin{sustechpipelinediagram}
    \sustechphasebox{inp}{at={(0,0)}}{Input}
    \sustechphasebox{pp}{right=of inp}{Preprocess}
    \sustechphasebox{feat}{right=of pp}{Features}
    \sustechphasebox{model}{right=of feat}{Model}
    \sustechphasebox{out}{right=of model}{Output}

    \sustechconnect{inp}{pp}
    \sustechconnect{pp}{feat}
    \sustechconnectwith{sustech highlight arrow}{feat}{model}
    \sustechconnect{model}{out}
    \sustechconnecttowith{sustech feedback arrow}[out=180,in=-90][Retrain]{out}{feat}
  \end{sustechpipelinediagram}
  \caption{Full-feature benchmark pipeline diagram}
\end{figure}

% ------------------------------------------------------------
\section{Conclusion}

The template compiles all subsystems in a single pass sequence.
See \textcite{lamport1994latex} and \parencite{shannon1948mathematical}
for background on the typesetting and information theory foundations.
Refer to \cref{tab:perf-methods} for the performance numbers.

\sustechprintbibliography

\end{document}
